\section{Optimización y Entrenamiento}

\begin{frame}{Función de Pérdida}
  \begin{columns}
    \begin{column}{0.5\textwidth}
      \textbf{Objetivo de entrenamiento:}
      \[
        \mathcal{L} = (1-\lambda)\mathcal{L}_1 + \lambda\,\mathcal{L}_{\text{D-SSIM}}, \quad \lambda = 0.2
      \]
      \begin{itemize}
        \item $\mathcal{L}_1$: precisión pixel a pixel
        \item $\mathcal{L}_{\text{D-SSIM}}$: preserva detalle estructural
      \end{itemize}
    \end{column}
    \begin{column}{0.5\textwidth}
      \textbf{Optimizador Adam} con tasas por grupo:
      \small
      \begin{tabular}{lc}
        \toprule
        Parámetro & LR inicial \\
        \midrule
        Posición $\boldsymbol{\mu}$ & $1.6\times10^{-4}$ \\
        Escala $\mathbf{s}$ & $5\times10^{-3}$ \\
        Rotación $\mathbf{q}$ & $10^{-3}$ \\
        Opacidad $o$ & $5\times10^{-2}$ \\
        SH grado 0 & $2.5\times10^{-3}$ \\
        SH grados 1--3 & $2.5\times10^{-4}$ \\
        \bottomrule
      \end{tabular}
    \end{column}
  \end{columns}
\end{frame}

\begin{frame}{Control Adaptativo de Densidad}
  \begin{columns}
    \begin{column}{0.5\textwidth}
      \textbf{Criterio de activación} (cada 100 iter, entre iter 500--15000):
      \[
        \bar{\nabla}_i = \frac{1}{K}\sum_k \|\nabla_{\boldsymbol{\mu}'_i}\mathcal{L}\|, \quad \bar{\nabla}_i > \tau_{\text{grad}} = 2\times10^{-4}
      \]
      \vspace{0.5em}
      \begin{block}{Clonar}
        Gaussiana \textbf{pequeña} ($\max s_k < \tau_{\text{size}}$)\\
        $\to$ copia exacta en la misma posición
      \end{block}
      \begin{block}{Dividir}
        Gaussiana \textbf{grande} ($\max s_k \geq \tau_{\text{size}}$)\\
        $\to$ dos nuevas a escala $s_k / 1.6$
      \end{block}
    \end{column}
    \begin{column}{0.5\textwidth}
      \begin{tikzpicture}[scale=0.85]
        % Gran gaussiana
        \node[draw=MidnightBlue, fill=MidnightBlue!15, ellipse, minimum width=2.5cm, minimum height=1.2cm] (big) at (0,1.5) {\small Grande};
        % Dos pequeñas
        \node[draw=ForestGreen, fill=ForestGreen!15, ellipse, minimum width=1.1cm, minimum height=0.6cm] (s1) at (-1.2,-0.2) {\tiny split};
        \node[draw=ForestGreen, fill=ForestGreen!15, ellipse, minimum width=1.1cm, minimum height=0.6cm] (s2) at (1.2,-0.2) {\tiny split};
        \draw[->, thick] (big) -- (s1);
        \draw[->, thick] (big) -- (s2);
        \node[below] at (0,-0.8) {\small Dividir};
        % Clonar
        \node[draw=accentorange, fill=accentorange!15, ellipse, minimum width=0.9cm, minimum height=0.55cm] (sm) at (-3.5,1.5) {\tiny peq};
        \node[draw=accentorange, fill=accentorange!15, ellipse, minimum width=0.9cm, minimum height=0.55cm] (cl) at (-3.5,0.3) {\tiny clon};
        \draw[->, thick] (sm) -- (cl);
        \node[below] at (-3.5,-0.2) {\small Clonar};
      \end{tikzpicture}
    \end{column}
  \end{columns}
\end{frame}

\begin{frame}{Poda y Reset de Opacidad}
  \begin{columns}
    \begin{column}{0.5\textwidth}
      \textbf{Poda} (misma frecuencia que densificación):
      \vspace{0.3em}
      \begin{itemize}
        \item Eliminar si $o_i < 0.005$
        \item Eliminar si tamaño en mundo $> \tau_w$
        \item Eliminar si radio 2D $> \tau_v$
      \end{itemize}
    \end{column}
    \begin{column}{0.5\textwidth}
      \begin{alertblock}{Reset periódico de opacidades (cada 3000 iter)}
        \begin{itemize}
          \item Todas las opacidades $\leftarrow \approx 0.01$
          \item Las Gaussianas activas se recuperan rápido
          \item Las redundantes son podadas
        \end{itemize}
      \end{alertblock}
    \end{column}
  \end{columns}
  \vspace{0.5em}
  \centering
  \textit{Este mecanismo evita la acumulación de Gaussianas fantasma sin utilidad.}
\end{frame}
